The three case studies presented in the previous section provide complementary perspectives on the EpiMDE platform's capabilities, limitations, and applicability across different epidemiological modeling scenarios. This discussion synthesizes the key insights from these studies, addressing questions of generality, usability, and platform maturity.

\subsection{Key Insights and Comparison between HIV and COVID-19 Modeling}

The HIV case study provides several important insights that distinguish it from the COVID-19 case study and demonstrate EpiMDE's broader applicability. First, \textbf{behavioral stratification} differs fundamentally from demographic stratification. While COVID-19 models typically stratify by age and comorbidity (biologically determined characteristics), HIV modeling requires stratification by sexual behavior, which involves complex social and behavioral factors. EpiMDE's Group and Product mechanisms handle both types of stratification uniformly, demonstrating the generality of the approach.

Second, \textbf{heterogeneous contact patterns} are more critical in HIV than in COVID-19. HIV transmission rates vary dramatically between population groups (e.g., homosexual men have higher transmission rates than heterosexual populations), and cross-group transmission is asymmetric (e.g., transmission from homosexual men to women occurs, but not vice versa within the same epidemiological context). EpiMDE's StratumSpecificRate feature was essential for capturing this heterogeneity, allowing different transmission rates for each ContactFlow depending on both the source and target strata.

Third, \textbf{chronic disease dynamics} differ from acute infectious diseases. HIV involves long-term treatment states (Treated with ART) and progressive disease stages (Living with AIDS) that do not appear in typical acute disease models like COVID-19. The model also incorporates demographic recruitment (ExternalSources) representing ongoing entry into sexually active populations, which is less relevant for short-term epidemic modeling but critical for endemic diseases studied over decades.

Fourth, the HIV case study validated EpiMDE's \textbf{usability for non-MDE experts}. The model construction required only epidemiological domain knowledge, without requiring understanding of metamodel extensions, wrapper abstractions, or advanced MDE concepts. This confirms that EpiMDE successfully lowers the barrier to entry for epidemiologists while maintaining sufficient expressiveness for complex models.

Finally, the HIV modeling experience revealed that EpiMDE is particularly well-suited for \textbf{diseases with established, stable dynamics}. Unlike COVID-19, which required frequent model updates as new variants and interventions emerged, HIV transmission dynamics are relatively stable. This suggests that EpiMDE is ideal for modeling endemic diseases and established epidemics, while the extensible metamodel may be more appropriate for rapidly evolving diseases requiring frequent structural modifications.

\subsection{Key Insights on the Rapid Prototyping workflow}
This example is not intended as a validated workflow or method, but as a stress test of the platform’s extensibility. Our illustrative example shows that EpiMDE is a viable platform on which to develop  domain-specific modeling workflows, such as interactive rapid prototyping.
This is due to its extensible architecture and appropriate metamodel, using the infrastructure of the Eclipse Modeling Framework. Overall, we are able to confirm the core assumption of MDE: that general-purpose model engineering platforms, when tailored appropriately, can provide significant value in high-stakes domains like epidemiology. 

The example further illuminates some actionable limitations in EpiMDE.
One challenge lies in preparing existing models for use with EpiMDE. 
At present, users must manually reformat their models to conform to the EpiMDE metamodel. 
This is a nontrivial and error-prone task, that can be a significant barrier to using a domain-specific workflow like rapid prototyping. 
However, as EpiMDE is built on the Eclipse Modeling Framework, 
it is well-suited to  leverage model transformations to automate this process. 


Another challenge we identified lies in the platform’s limited support for semantic distinctions across several steps of the workflow.
Currently, our tooling identifies variation points through exact identifier string matching. This  fails to capture semantically equivalent elements, e.g., compartments labeled $S$ and $Susceptible$ are treated as distinct, despite referring to the same epidemiological concept.
Similarly, its clustering mechanisms are based on structural co-occurrence, without regard to the conceptual roles of model elements. 
For example, we do not identify that each input model represents infectious populations differently, and we depend on interactive input for the user to add the appropriate logical relationships.
The detection of logical dependencies between features is also syntactic, flagging structural inconsistencies (e.g., disconnected flows) but missing domain-specific conflicts, such as including both asymptomatic transmission and severe hospitalization without intermediate transitions.
These limitations point to an opportunity to enrich the platform’s semantic processing capabilities.
Future work can explore tighter integration of EpiMDE with Natural Language Processing techniques and Large Language Models, to enable more flexible variation point matching, more coherent clustering, and the identification of plausible feature combinations based on domain knowledge.
This would reduce the manual burden on users and improve the quality and relevance of generated model prototypes.

%\textbf{Limitations in Variation Point Identification: }In step 1 of our approach,  variation point IDs for model elements rely on exact matches of element properties. While this approach simplifies comparison, it misses semantically equivalent elements. For example, variation points for compartments are solely based on compartment names. This could treat compartments named $S$ and $Susceptible$ as distinct variation points, even though they likely represent the same epidemiological concept.
%To address this limitation, future work could use the descriptions that often accompany compartments and apply natural language processing (NLP) techniques to recognize similarities in meaning rather than exact term matching.

%\textbf{Lack of Semantic Refinement in Clustering: }The clusters created using Formal Concept Analysis (FCA) in step 2 are purely structural and only based on the co-occurrence of variation points across models. However, considering the underlying meaning of the elements could avoid clusters with semantically unrelated elements.
%Although we offer an interactive step for the epidemiologist to refine the clusters, a future improvement could design an additional step of semantic refinement before presenting the clusters to the user. This refinement could consider the names, descriptions, and context of the variation points to reorganize clusters to be more coherent.

%\textbf{Limited Detection of Semantic Dependencies: }The automatically identified logical dependencies between features in step 3 are limited to syntactic relationships, such as conflicting flows or missing source and target compartments. While this captures many important constraints, epidemiologists must still manually add semantic or domain-specific dependencies. Future work could streamline users' work by integrating natural language processing techniques to detect potential semantic conflicts or unlikely feature combinations. For instance, if one feature models asymptomatic transmission while another models hospitalization due to severe symptoms, it may be unlikely or illogical for an epidemiologist to include both in the same prototype model. 

\subsection{Key Insights on the LLM-based Automated Model Generation}

\paragraph{Shared Challenges} Both models struggled with \textbf{indexing errors}—incorrect 0-based indices in flow references (e.g., \texttt{target="//@compartments.X"}). The EpiMDE metamodel requires precise numerical tracking across dozens of elements, a task that strains current LLMs' symbolic reference capabilities despite their strong semantic understanding. This represents a fundamental limitation: LLMs excel at pattern matching and natural language understanding but struggle with precise bookkeeping of numerical relationships across large structured outputs.

\paragraph{Stage 3 Robustness} Both models achieved 100\% execution success for simulation script generation (Stage 3) across both case studies. The two-phase skeleton approach—pre-writing universal components (imports, loop structure, plotting) and narrowly scoping LLM responsibilities to ODE transcription—successfully isolated code generation from the error-prone aspects of XML generation. This validates task decomposition as an effective design pattern for LLM-based scientific code generation.

\paragraph{Implications} The cross-model comparison demonstrates that no single LLM currently provides uniformly superior performance. However, several insights emerge. First, \textbf{frontier model access is critical}: Gemini's perfect COVID-19 generation versus GPT-4o's 65\% recall strongly suggests that complex epidemiological model generation requires state-of-the-art capabilities. As models grow in complexity, access to cutting-edge LLMs transitions from optional to necessary. Second, \textbf{error mode complementarity}: Gemini's perfect recall with imperfect precision (hallucinations) is arguably preferable to GPT-4o's imperfect recall (omissions), as false positives are easier to detect and correct in post-processing than false negatives. Validation layers can flag unexpected attributes, but cannot identify missing flows without ground truth comparison. Third, \textbf{indexing requires targeted intervention}: Since both models struggle with numerical reference tracking, prompt engineering alone may be insufficient. Post-processing validation (checking reference consistency) or intermediate index mapping steps (explicitly documenting element positions before XML generation) are likely necessary complements.

\paragraph{Limitations and Challenges} While the pipeline demonstrates substantial promise, several limitations warrant discussion. First, \textbf{validation overhead}: LLM outputs require careful validation, particularly for complex models where omissions or hallucinations can silently corrupt model semantics. The pipeline is best viewed as an assistive technology accelerating modeling rather than a fully autonomous system. Second, \textbf{prompt brittleness}: Performance depends heavily on prompt structure and clarity. The COVID-19 flow omission pattern in GPT-4o suggests that even clear user specifications can be ambiguous to LLMs if prompt instructions for translating descriptions into XML are insufficiently explicit. Third, \textbf{semantic verification gaps}: The pipeline enforces structural constraints (valid XML, correct element nesting) but cannot detect epidemiologically implausible parameter values or logically invalid model structures (e.g., transmission without contact, recovery before infection). Domain expertise remains essential for semantic validation. Fourth, \textbf{scalability uncertainty}: While Gemini handled the 63-element COVID-19 model perfectly, we have not yet tested models with hundreds of compartments or dozens of stratification dimensions. Performance degradation at extreme scales remains possible.

\subsection{Key Insights on the EpiMDE Modeling Platform}

\paragraph{Accessibility and Immediate Usability}

A central design goal of EpiMDELite was to provide immediate usability for epidemiologists without requiring MDE expertise. The HIV and COVID-19 case studies validate this goal. Both studies successfully replicated published mathematical models, producing simulation results identical to the original publications. The HIV case study, in particular, demonstrated that epidemiologists can leverage EpiMDELite's stratification features (Groups, Products, StratumSpecificRate) without understanding the underlying metamodel mechanics. The ability to define a single Group for "SexualBehavior" and automatically expand compartments across three demographic strata represents a significant reduction in manual effort compared to traditional tools like IDM-CMS or Matlab, where modelers must manually write equations for each stratified compartment.

However, the rapid prototyping case study revealed a current limitation in accessibility: preparing existing models for use with EpiMDE requires manual reformatting to conform to the metamodel. This creates a barrier to entry when epidemiologists want to import legacy models or models from collaborators. While model transformation techniques can address this limitation, the current manual process highlights that accessibility is not yet complete across all platform workflows.

\paragraph{Stratification as a First-Class Concept}

Population stratification emerged as a critical differentiator across the case studies. The HIV model required behavioral stratification (sexual behavior), while the COVID-19 model required demographic stratification (age, comorbidity). EpiMDELite's unified approach to stratification through Groups and Products handled both cases uniformly, despite their conceptual differences. This contrasts with traditional approaches where demographic stratification might be handled through explicit compartment multiplication, while behavioral stratification requires careful consideration of contact patterns.

The COVID-19 study further highlighted the value of explicit stratification representation. The original Tuite et al. model stratified by 16 age groups and comorbidity status, but this stratification was "invisible" in the graphical representation and only appeared in the equations. EpiMDE makes stratification explicit in the model structure, improving transparency and maintainability. This explicitness also facilitated model validation, as the stratification structure was immediately apparent rather than buried in mathematical notation.

\paragraph{Disease Characteristics and Metamodel Selection}

The case studies revealed important insights about when to use EpiMDELite versus the extensible metamodel. EpiMDELite proved well-suited for diseases with established, stable dynamics (HIV) where the core model structure is unlikely to change frequently. For such diseases, the simplified metamodel's constraints are not limiting, and its accessibility benefits dominate.

In contrast, the COVID-19 study demonstrated scenarios where the extensible metamodel's additional capabilities become valuable. COVID-19's rapid evolution (new variants, changing interventions) requires frequent model restructuring. The extensible metamodel's wrapper-based abstractions and compositional features (e.g., Groups within Products, Link compartments) enable modular model evolution. For instance, adding isolation status as a new dimension or incorporating hospitalization dynamics could be done through metamodel extensions without restructuring the entire model.

This suggests a maturity progression: epidemiologists can begin with EpiMDELite for initial model development and validation, then transition to the extensible metamodel when their modeling needs require more sophisticated structural manipulation.

\paragraph{Reproducibility and Model Sharing}

The first three case studies emphasized reproducibility challenges in current epidemiological practice. The COVID-19 study explicitly noted that critical model information was only available in the source code repository, not the published manuscript. The rapid prototyping study addressed this through systematic model reuse, but revealed that semantic matching (e.g., recognizing that "S" and "Susceptible" refer to the same concept) remains a manual task.

EpiMDE improves reproducibility through several mechanisms. First, graphical model representations serve as executable documentation, capturing both structure and semantics in a single artifact. Second, automatic equation generation eliminates transcription errors that plague manual equation implementation. Third, the platform's model differencing capabilities (demonstrated in the COVID-19 validation) enable systematic comparison between models, facilitating validation and knowledge transfer between research groups.

However, the rapid prototyping study identified an important gap: while EpiMDE facilitates reuse within the platform, importing models from external sources (e.g., published papers, other tools) remains manual. Future integration with model transformation frameworks could address this by providing automated translation from common epidemiological modeling formats.

\paragraph{Validation and Correctness}

Model validation emerged as both a strength and an opportunity for improvement. The HIV and COVID-19 studies demonstrated that EpiMDE-generated simulations match published results exactly, validating the platform's correctness. The COVID-19 study's validation approach (comparing derivatives over randomly generated population distributions) provides a systematic method for verifying model equivalence.

However, the rapid prototyping study revealed limitations in semantic validation. The platform detects structural inconsistencies (e.g., flows with missing source/target compartments) but cannot identify domain-specific logical errors (e.g., modeling asymptomatic transmission alongside severe hospitalization without intermediate states). While metamodel constraints prevent many errors at design time, richer domain-specific validation rules could further improve model quality.

\paragraph{Platform Extensibility and Workflow Support}

The rapid prototyping case study demonstrates that EpiMDE's infrastructure (built on EMF and Eclipse) supports the development of domain-specific workflows beyond basic modeling. The four-step rapid prototyping workflow (variation point extraction, feature identification, relationship identification, prototype generation) was implemented as an extension of the platform without modifying the core metamodel or tooling.

This extensibility validates the MDE approach's core assumption: that general-purpose model engineering platforms, when appropriately tailored, can provide significant value in specialized domains. However, the study also revealed friction points. The reliance on exact string matching for variation point identification and purely structural clustering limits semantic reasoning. Future integration with natural language processing or large language models could enhance semantic capabilities while maintaining the platform's structural rigor.

\paragraph{Automation and AI-Assisted Modeling}

Rather than simplifying the modeling interface, the LLM pipeline eliminates the need to use it for initial model creation. This represents an important evolution in making MDE accessible to domain experts: instead of training users on metamodels and graphical editors, the system accepts specifications in formats already familiar to epidemiologists (textual descriptions with mathematical notation).

The multi-stage pipeline architecture (structural generation, value mapping, simulation synthesis) demonstrates that LLMs can serve as effective translators between domain-specific natural language and formal metamodel-compliant representations. The key insight is constraining LLM behavior through structured prompts and explicit validation stages rather than relying on end-to-end generation. This design acknowledges LLM limitations (potential for hallucination, difficulty with precise syntax) while leveraging their strengths (understanding natural language, pattern recognition in mathematical notation).

The LLM pipeline complements rather than replaces the graphical IDE. Models generated automatically can be visualized, validated, and refined graphically. This hybrid workflow accommodates different user preferences and modeling scenarios: textual specifications for reproducing published models, graphical editing for exploratory modeling and composition. The case study highlights that accessibility in MDE should be viewed as a spectrum of interaction modalities rather than a single "best" interface.

However, the LLM approach also introduces new challenges. Validation becomes more critical when models are generated automatically, as users may not immediately understand the resulting structure. The pipeline's reliance on careful prompt engineering means that robustness depends on anticipating potential failure modes. Future work should explore interactive validation workflows where LLMs explain their generated models to users, enabling collaborative refinement rather than one-shot generation.

\paragraph{Comparison to Traditional Approaches}

Across all case studies, EpiMDE demonstrated advantages over traditional epidemiological modeling tools. Compared to tools like STEM (complex, crash-prone interface) or IDM-CMS (hundreds of lines of manual equation writing), EpiMDE's graphical modeling and automatic equation generation significantly reduce cognitive load and development time. The HIV model, which would require separate equations for each of 15 stratified compartments (5 disease states × 3 behavioral groups) in traditional tools, was constructed using a single Product definition in EpiMDE.

However, traditional tools retain advantages in specific areas. Tools like Matlab offer greater flexibility for non-standard model structures, while R-based tools like EpiModel provide extensive statistical analysis capabilities that EpiMDE currently lacks. EpiMDE's value proposition is strongest when model structure, stratification, and reproducibility are primary concerns, rather than statistical inference or highly customized dynamics.

\paragraph{Implications for MDE in Scientific Computing}

The case studies provide broader lessons for applying MDE in scientific computing domains. First, accessibility cannot be assumed even when using domain-familiar concepts. The manual reformatting required for legacy models highlights that interoperability with existing ecosystems is critical for adoption. The LLM pipeline addresses this by accepting natural language specifications, demonstrating that AI-assisted transformation can bridge the gap between traditional scientific workflows and formal metamodeling.

Second, metamodel design involves fundamental trade-offs between simplicity and expressiveness. EpiMDELite's success in the HIV study and the extensible metamodel's value in COVID-19 modeling suggest that offering both options, with clear guidance on when to use each, may be more effective than seeking a single "optimal" metamodel. Third, semantic reasoning remains challenging in MDE approaches. While metamodels excel at capturing structure and enforcing structural constraints, domain-specific semantic validation often requires complementary techniques. The LLM pipeline demonstrates one such complementary approach, though it introduces new validation requirements.

Fourth, the LLM pipeline case study suggests that the future of MDE accessibility may lie in multi-modal interaction rather than finding the "right" interface. Different users and scenarios benefit from different interaction modes: textual specifications for reproducing literature, graphical editing for exploration, programmatic APIs for batch processing. Supporting seamless transitions between these modes may be more valuable than optimizing any single mode.

These insights suggest that successful MDE platforms for scientific domains should prioritize: (1) interoperability with existing tools and formats through automated transformation, (2) layered abstraction with clear migration paths between levels, (3) hybrid approaches combining structural metamodeling with semantic reasoning techniques from AI and natural language processing, and (4) multi-modal interaction supporting diverse user preferences and workflow requirements.